% Input fragment; use the preamble of the Chinese Volume II source.
% This is an ordinary chapter. Preserve existing special Chapter 5.5 counters.
\chapter{待定章节标题}\label{template:v2}
\begin{chaptersummary}
【作者待填写：约 100--250 字，连接已有结果、当前问题、
采用的结构或工具、后续离散用途和一个核心思想。】
\end{chaptersummary}
\studyroute{【明确略去的证明、代数或实现细节，并添加已核实的文献。】}
{【填写精确、通用的英文检索词】}
{【从当前定义或方程出发，指定可检查的展开任务。】}

\section{对象、约定与问题}
【沿用本卷的单位、度规号差和取向，说明本节使用的场、空间与假设。】

% Optional draft author input: Volume II has no authornote environment.
% studentinput below is a reader exercise, not an author placeholder.
\paragraph{潜在认知跳跃}
\textbf{潜在疑问：}【为什么此处需要引入新的对象或操作？】\par
\textbf{作者补充：}【作者待补充：80--300 字，说明动机与选择依据。】\par
\textbf{核心概念联系：}【已有结果】$\longrightarrow$【当前问题】。

\section{关键关系与计算}
【加入已确认的公式、边界项与解释，区分连续、半离散和全离散结论。】
\begin{askai}
【指定当前符号、假设、起点以及逐步推导或可复算的任务。】

进一步展开：【改变一个条件，检验结果的范围。】
\end{askai}

% Optional active-recall and implementation blocks.
\begin{studentinput}
【填写一个读者可以独立完成的具体任务。】

\textbf{你的答案：}\par\vspace{1cm}
\end{studentinput}
\begin{checkpoint}
【填写可检验核心概念区别的具体问题。】
\end{checkpoint}
\begin{engineernote}
【说明已确认的结构对自由度、接口或实现的影响及其适用条件。】
\end{engineernote}

\section{练习}
\begin{enumerate}
\item 【给出具有明确输入和检查目标的练习。】
\end{enumerate}
